\chapter{Conclusion}

\section{Summary and Key Findings}

This thesis investigated the classification of high-frequency partial discharge (PD) signals from medium-voltage overhead lines using traditional machine learning models under two independent feature selection regimes: Featurewiz and MLJAR. The study demonstrated that, despite substantial differences in feature construction, dimensionality, and selection strategy, both regimes consistently identified a stable core of informative descriptors spanning time-domain statistics, frequency-domain summaries, complexity measures, and phase-resolved patterns.

Comparative evaluation of machine learning models showed that ensemble-based and nonlinear classifiers achieved the strongest overall performance. Under the Featurewiz feature selection track, the Random Forest classifier attained the highest classification accuracy of $91.91\%$ with a corresponding F1-score of $71.84\%$, demonstrating competitive performance using a compact and interpretable feature set. In contrast, the MLJAR-based approach achieved its best performance using a Random Forest classifier with an accuracy of $91.83\%$, precision of $81.50\%$, recall of $68.94\%$, F1-score of $73.11\%$, and ROC-AUC of $90.63\%$.

While the MLJAR-based framework yielded a modest improvement in F1-score, this improvement was accompanied by a substantial increase in feature dimensionality, highlighting an important trade-off between classification performance and model complexity. The Featurewiz-based approach, by comparison, achieved near-equivalent accuracy with a significantly reduced feature set, offering advantages in computational efficiency and interpretability that are desirable for practical deployment.

Feature importance analysis further revealed that classical statistical indicators of impulsiveness and asymmetry remain diagnostically relevant, even when their contribution is implicitly embedded within higher-order interaction features. This observation clarifies that apparent differences in feature rankings across selection frameworks arise primarily from methodological representation rather than disagreement in the underlying physical characteristics of partial discharge phenomena.

Overall, the results confirm that robust PD classification can be achieved through multi-domain feature integration and model-agnostic feature selection strategies. The findings presented in this thesis provide a physically grounded and practically viable basis for the development of reliable partial discharge diagnostic systems suitable for real-world monitoring environments.


\section{Limitations of the Study}

A key limitation of this study arises from the exclusive reliance on long-term, multi-site field monitoring data acquired from operational overhead power lines, which are subject to inherently uncontrolled measurement conditions \parencite{Florkowski_2024, Sikorski_2025}. 
Field-acquired partial discharge (PD) signals are therefore unavoidably affected by environmental noise sources, including electromagnetic interference, weather-induced disturbances, and load-dependent fluctuations, which can obscure weak discharge signatures and reduce classification reliability \parencite{Orellana_2023_1, Florkowski_2024}. 
Further variability introduced by insulation aging, material heterogeneity, and temporal degradation processes complicates the separation of discharge-specific features from background interference \parencite{Seitz_2024_1, Sikorski_2025}. 
Although such real-world conditions enhance ecological validity and reflect practical deployment scenarios, they limit the isolation of fundamental PD mechanisms under high signal-to-noise clarity \parencite{Florkowski_2024, Long_2024}.

The absence of complementary experimental data acquired under controlled laboratory conditions restricts the evaluation of intrinsic feature sensitivity and noise-independent discriminative capability \parencite{Aldosari_2023, Long_2024}. 
Laboratory-based PD measurements, characterized by well-defined defect geometries, stable operating conditions, and high signal-to-noise ratios, provide an essential reference for validating feature relevance under idealized conditions \parencite{Aldosari_2023, Long_2024}. 
As a result, the generalizability of the proposed framework across the full spectrum of PD operating conditions cannot be fully established using field data alone \parencite{Seitz_2024_1, Florkowski_2024}. 
Future work should therefore integrate controlled experimental datasets with real-world monitoring data to enable systematic benchmarking between noise-free and operational environments and to strengthen the physical interpretability of PD classification models \parencite{Long_2024, Sikorski_2025}.

A further limitation concerns the geographic and climatic context of the training data \parencite{Florkowski_2024, Sikorski_2025}. 
The dataset utilized in this study was acquired from monitoring stations located in Czechia and Slovakia, regions characterized by temperate continental climates with distinct four-season annual cycles \parencite{Florkowski_2024, Sikorski_2025}. 
Such climatic variability introduces environmental conditions including temperature fluctuations, humidity patterns, precipitation regimes, and seasonal insulation stress variations that may not be representative of regions with different climate characteristics \parencite{Orellana_2023_1, Azam_2024}.

Specifically, four-season climate dynamics may influence PD activity through temperature-dependent insulation behavior, humidity-driven surface contamination effects, and seasonally varying electromagnetic interference patterns \parencite{Orellana_2023_1, Florkowski_2024}. 
The transferability of models trained on this data to regions with fundamentally different climate regimes, such as tropical or equatorial environments with only wet and dry seasons, remains unvalidated \parencite{Azam_2024, Florkowski_2024}. 
For example, deployment in Malaysia may encounter distinct PD signal characteristics, environmental noise profiles, and insulation stress mechanisms that could affect model performance \parencite{Azam_2024, Florkowski_2024}. 
This limitation underscores the importance of region-specific validation and potential domain adaptation strategies when deploying PD classification models across diverse geographic and climatic contexts \parencite{Seitz_2024_1, Friebe_2024}.

The presence of precision--recall performance gaps, where even the best-performing models achieved recall values ranging from 67.09\% to 72.32\%, indicates that a proportion of fault instances remain undetected despite improved sensitivity to the minority class \parencite{Zemouri_2023_1, Aldosari_2023}. 
This constraint motivates the adoption of threshold calibration, cost-sensitive learning strategies, and class rebalancing techniques in scenarios where false-negative risk is operationally critical \parencite{Zemouri_2023_1, Adhikari_2024}.

\section{Transition to Future Work}

External validation across additional assets, monitoring sites, and acquisition modalities constitutes a primary requirement for establishing the broader transferability and generalizability of the developed PD classification models beyond the present dataset \parencite{Seitz_2024_1, Sikorski_2025}. 
Such validation should explicitly include datasets collected from diverse geographic regions and climatic regimes to assess the impact of environmental variability on model performance, directly addressing the identified limitation associated with four-season versus two-season climate contexts \parencite{Azam_2024, Florkowski_2024}.

Differences in noise statistics, electromagnetic coupling paths, sensor configurations, and environmental interference patterns across deployment contexts further motivate the investigation of domain adaptation and robustness-oriented learning strategies \parencite{Friebe_2024, Florkowski_2024}. 
Techniques such as transfer learning, domain adversarial training, and environment-aware data augmentation could be explored to enhance model robustness under heterogeneous operational and climatic conditions \parencite{Seitz_2024_1, Azam_2024}.

In parallel, future work should focus on explicitly addressing class imbalance and asymmetric operational risk through calibrated decision thresholds, cost-sensitive learning objectives, and advanced resampling strategies to improve minority fault detection sensitivity \parencite{Zemouri_2023_1, Adhikari_2024}. 
Additionally, data augmentation and advanced denoising approaches are widely recognized as practical mechanisms for improving robustness to interference and enhancing generalization when operating on noisy field-acquired PD signals \parencite{Bao_2024_1, Friebe_2024}.

Finally, principled feature set reduction combined with interpretability analysis and uncertainty quantification should be pursued to support deployable monitoring workflows that balance classification performance, computational efficiency, and engineering trust \parencite{Vera_2023_1, Florkowski_2024}. 
The integration of explainable artificial intelligence techniques with the developed models would further enhance their suitability for industrial deployment scenarios where transparent diagnostics and confidence-aware decision-making are critical for asset maintenance planning \parencite{Vera_2023_1, Seitz_2024_1}.

Beyond the evaluated classifiers, future studies may explore a broader range of machine learning and deep learning architectures, including ensemble hybrids, attention-based neural networks, and transformer-inspired models, to capture more complex temporal and spectral dependencies in PD signals \parencite{Zemouri_2023_1, Seitz_2024_1}. 
Such models have demonstrated improved representation learning capabilities in noisy and non-stationary signal environments, making them promising candidates for real-world PD monitoring applications \parencite{Florkowski_2024, Sikorski_2025}.

In addition, the adoption of advanced denoising techniques, such as wavelet-domain adaptive filtering, empirical mode decomposition variants, and learning-based denoising autoencoders, could further improve signal quality prior to feature extraction and classification \parencite{Bao_2024_1, Friebe_2024}. 
These approaches enable the suppression of non-stationary interference while preserving discharge-relevant signal components, which is critical for enhancing classification robustness under field-acquired noise conditions \parencite{Orellana_2023_1, Florkowski_2024}.